\documentclass{article}
\usepackage{import}
\import{../../template/}{article-template}
\graphicspath{{../../images/}}
\addbibresource{../../template/library.bib}

\title{Haskell}
\author{PENG}
\date{Created: 20260913, Version: 20260913}

\begin{document}
\maketitle
\tableofcontents
\section{Introduction}
\url{https://www.engineering.upenn.edu/~cis1940/spring13/}

Haskell is a \textit{lazy}, \textit{functional} programming language created in the late 1980s.

\subsection{Functional}
\begin{definition}[functional]~
    \begin{itemize}
        \item functions are values
        \item evaluating expressions rather than executing instructions
    \end{itemize}
\end{definition}

\subsection{Pure}
\begin{definition}[referentially transparent]~
    \begin{itemize}
        \item everything (var, data structures) is immutable
        \item expressions have no side effects
        \item same function and arguments = same output
    \end{itemize}
\end{definition}

benefits:
\begin{itemize}
    \item replace equals by equals
    \item parallelism
\end{itemize}

\subsection{Lazy}
\begin{definition}[lazy]
    Haskell expressions are not evaluated until their results are actually needed.
\end{definition}

consequences:
\begin{itemize}
    \item define a control structure by defining a function
    \item define infinite data structures
    \item time and space usage become more complicated
\end{itemize}

\subsection{Statically Typed}
\begin{definition}[statically typed]
    Each Haskell expression has a type, and types are all checked at compile-time.
\end{definition}

\subsection{Themes}
\begin{itemize}
    \item static type system
    \item abstraction principle:
          \textit{every idea, algorithm, and piece of data should occur exactly once in your code}
    \item wholemeal programming:
          \textit{first solve a more general problem, then transform it into more specialized ones}
\end{itemize}

\subsection{Declarations and variables}
\begin{definition}[=]
    = denotes definition such that variables are not mutable
\end{definition}

\begin{example}
    \begin{verbatim}
x :: Int
x = 3
x = 4 -- error, multiple declaration of x
y :: Int
y = y + 1 -- <<loop>>
\end{verbatim}
\end{example}

\subsection{Basic Types}
\begin{itemize}
    \item \verb|Int| : at least up to $\pm 2^{29}$
    \item \verb|Integer|
    \item \verb|Double|
    \item \verb|Float|
    \item \verb|Bool|
    \item \verb|Char|
    \item \verb|String|
\end{itemize}

\subsection{Arithmetic}
\begin{example}~
    \begin{verbatim}
mod 18 5
18 `mod` 5
2 / 2 -- floating-point division
div 2 2 -- integer division
2 `div` 2
\end{verbatim}
\end{example}

\subsection{Boolean Logic}
\begin{definition}[boolean logic]
    \verb|&&|, \verb~||~, \verb|not|
\end{definition}

\begin{definition}[equality and order]~
    \begin{itemize}
        \item \verb|==| or \verb|/=|
        \item \verb|>|, \verb|<|, \verb|>=|, \verb|<=|
    \end{itemize}
\end{definition}

\begin{definition}[if-expression]
    \verb|if b then t else f|
\end{definition}

\begin{remark}
    For if-expression, the \verb|else| part is required.
\end{remark}

\subsection{Defining Basic Functions}


\printbibliography
\end{document}
